% !TeX program = xelatex
\documentclass[a4paper,12pt]{article}

% ===== 1. KHỔ GIẤY, LỀ TRANG, FONT =====
\usepackage[left=3cm,right=2cm,top=2cm,bottom=2cm,footskip=1cm]{geometry}
\usepackage{fontspec}

\IfFontExistsTF{Times New Roman}{
  \setmainfont{Times New Roman}
}{
  \IfFileExists{fonts/times.ttf}{
    \setmainfont{times.ttf}[
      Path=fonts/,
      BoldFont=timesbd.ttf,
      ItalicFont=timesi.ttf,
      BoldItalicFont=timesbi.ttf]
  }{
    \IfFontExistsTF{TeX Gyre Termes}{
      \setmainfont{TeX Gyre Termes}
      \PackageWarningNoLine{MauBaoCao}{Times New Roman unavailable; using TeX Gyre Termes}
    }{
      \setmainfont{DejaVu Serif}
      \PackageWarningNoLine{MauBaoCao}{Times New Roman unavailable; using DejaVu Serif}
    }
  }
}
\usepackage{indentfirst}
\usepackage{microtype}
\usepackage{graphicx}
\usepackage{float}
\usepackage{tikz}
\usepackage{array,tabularx,booktabs,longtable}
\usepackage{enumitem}
\usepackage{titlesec}
\usepackage{tocloft}
\usepackage{caption}
\usepackage{xurl}
\usepackage[unicode,hidelinks]{hyperref}
\urlstyle{same}

\renewcommand{\normalsize}{\fontsize{13}{19.5}\selectfont}
\AtBeginDocument{\normalsize}
\setlength{\parindent}{0.5cm}
\setlength{\parskip}{4pt}
\setlength{\emergencystretch}{2em}
\raggedbottom
\setlist{topsep=4pt,itemsep=2pt,parsep=0pt,leftmargin=*}

% ===== 2. BA CẤP TIÊU ĐỀ: 1 / 1.1 / 1.1.1 =====
\setcounter{secnumdepth}{3}
\setcounter{tocdepth}{3}
\renewcommand{\thesection}{\arabic{section}}
\renewcommand{\thesubsection}{\thesection.\arabic{subsection}}
\renewcommand{\thesubsubsection}{\thesubsection.\arabic{subsubsection}}
\titleformat{\section}[hang]
  {\normalfont\fontsize{16}{20}\selectfont\bfseries}
  {\thesection.}{0.6em}{}
\titleformat{\subsection}[hang]
  {\normalfont\fontsize{14}{18}\selectfont\bfseries}
  {\thesubsection}{0.6em}{}
\titleformat{\subsubsection}[hang]
  {\normalfont\fontsize{13}{17}\selectfont\bfseries}
  {\thesubsubsection}{0.6em}{}
\titlespacing*{\section}{0pt}{16pt}{8pt}
\titlespacing*{\subsection}{0pt}{12pt}{6pt}
\titlespacing*{\subsubsection}{0pt}{8pt}{4pt}

% ===== 3. MỤC LỤC VÀ SỐ TRANG =====
\renewcommand{\contentsname}{MỤC LỤC}
\renewcommand{\refname}{Tài liệu tham khảo}
\renewcommand{\cfttoctitlefont}{\fontsize{16}{20}\selectfont\bfseries}
\setlength{\cftbeforetoctitleskip}{0pt}
\setlength{\cftaftertoctitleskip}{18pt}
\renewcommand{\cftsecfont}{\normalfont\bfseries}
\renewcommand{\cftsecpagefont}{\normalfont\bfseries}
\renewcommand{\cftsubsecfont}{\normalfont}
\renewcommand{\cftsubsubsecfont}{\normalfont}
\renewcommand{\cftsecleader}{\cftdotfill{\cftdotsep}}
\cftsetindents{section}{0em}{2em}
\cftsetindents{subsection}{2em}{3em}
\cftsetindents{subsubsection}{5em}{4em}
\setlength{\cftbeforesecskip}{6pt}
\makeatletter
\def\ps@plain{%
  \let\@oddhead\@empty\let\@evenhead\@empty
  \def\@oddfoot{\hfil{\normalfont\fontsize{11}{13}\selectfont\thepage}\hfil}%
  \let\@evenfoot\@oddfoot}
\makeatother
\pagestyle{plain}

% ===== 4. BẢNG, HÌNH, CHÚ THÍCH =====
\renewcommand{\tablename}{Bảng}
\renewcommand{\figurename}{Hình}
\DeclareCaptionFont{reportcaption}{\fontsize{12}{15}\selectfont}
\captionsetup{font=reportcaption,labelfont=bf,labelsep=period,
  justification=centering,skip=6pt}
\newcolumntype{Y}{>{\raggedright\arraybackslash}X}
\newcommand{\bangfont}{\fontsize{12}{15}\selectfont}
\newcommand{\code}[1]{\texttt{#1}}

% ===== 5. THÔNG TIN TRANG BÌA =====
\newcommand{\DonViChuQuan}{ĐẠI HỌC QUỐC GIA HÀ NỘI}
\newcommand{\TenTruong}{TRƯỜNG ĐẠI HỌC CÔNG NGHỆ}
\newcommand{\LoaiBaoCao}{BÁO CÁO TIẾN ĐỘ DỰ ÁN 2}
\newcommand{\TenMonHoc}{ĐIỆN TOÁN ĐÁM MÂY}
\newcommand{\TenDeTai}{SatiBot -- Quản lý và tương tác với các tài nguyên AWS bằng ngôn ngữ tự nhiên}
\newcommand{\GiangVien}{TS. Hoàng Xuân Tùng}
\newcommand{\SoNhom}{Nhóm 10}
\newcommand{\LopMonHoc}{INT3319 1}
\newcommand{\DiaDiemThoiGian}{Hà Nội, Tháng 10/2026}
\hypersetup{pdftitle={BÁO CÁO TIẾN ĐỘ DỰ ÁN 2 - SATIBOT},pdfauthor={Nhóm 10}}

\begin{document}

% ===== 6. TRANG BÌA =====
\begin{titlepage}
\thispagestyle{empty}
\begin{tikzpicture}[remember picture,overlay]
  \draw[line width=2pt]
    ([xshift=3cm,yshift=-2cm]current page.north west)
    rectangle ([xshift=-2cm,yshift=2cm]current page.south east);
  \draw[line width=0.5pt]
    ([xshift=3.12cm,yshift=-2.12cm]current page.north west)
    rectangle ([xshift=-2.12cm,yshift=2.12cm]current page.south east);
\end{tikzpicture}\par
\centering
\setlength{\parindent}{0pt}
\setlength{\parskip}{0pt}
\vspace*{0.8cm}
{\fontsize{14}{18}\selectfont\bfseries\DonViChuQuan\par}
\vspace{0.15cm}
{\fontsize{15}{20}\selectfont\bfseries\TenTruong\par}
\vspace{0.25cm}
\rule{5cm}{0.4pt}\par
\vspace{0.5cm}

\IfFileExists{logo.png}{
  \includegraphics[width=3.8cm,height=3.8cm,keepaspectratio]{logo.png}
}{
  \IfFileExists{Report progress 2/logo.png}{
    \includegraphics[width=3.8cm,height=3.8cm,keepaspectratio]{Report progress 2/logo.png}
  }{
    \IfFileExists{Report progress 1/logo.png}{
      \includegraphics[width=3.8cm,height=3.8cm,keepaspectratio]{Report progress 1/logo.png}
    }{
      \fbox{\parbox[c][3.4cm][c]{3.4cm}{\centering
        \fontsize{11}{14}\selectfont LOGO TRƯỜNG}}
    }
  }
}\par
\vspace{0.7cm}
{\fontsize{22}{27}\selectfont\bfseries\LoaiBaoCao\par}
\vspace{0.35cm}
\begin{minipage}{0.9\textwidth}\centering
  {\fontsize{15}{20}\selectfont\bfseries MÔN HỌC: \TenMonHoc\par}
\end{minipage}\par
\vspace{0.6cm}
{\fontsize{13}{17}\selectfont\bfseries ĐỀ TÀI\par}
\vspace{0.12cm}
\begin{minipage}{0.9\textwidth}\centering
  {\fontsize{16}{21}\selectfont\bfseries\TenDeTai\par}
\end{minipage}\par
\vspace{0.65cm}
{\fontsize{13}{19.5}\selectfont
\renewcommand{\arraystretch}{1.1}
\begin{tabular}{@{}>{\raggedright\arraybackslash\bfseries}p{5.2cm}>{\raggedright\arraybackslash}p{8cm}@{}}
  Giảng viên hướng dẫn: & \GiangVien\\
  Nhóm thực hiện: & \SoNhom\\
  Sinh viên thực hiện: & Cao Anh Minh \textbf{(Trưởng nhóm)}\\
                      & Trịnh Thiên Quang\\
                      & Phạm Đức Hùng\\
                      & Nguyễn Đức Huy\\
  Lớp môn học: &  \LopMonHoc
\end{tabular}\par}
\vfill
{\fontsize{13}{17}\selectfont\bfseries\DiaDiemThoiGian\par}
\vspace*{0.5cm}
\end{titlepage}

% ===== 7. MỤC LỤC TỰ ĐỘNG =====
\pagenumbering{roman}
\tableofcontents
\clearpage
\pagenumbering{arabic}

% ===== 8. NỘI DUNG BÁO CÁO =====

\section{TỔNG QUAN VỀ DỰ ÁN \textit{(không thay đổi)}}

\subsection{Giới thiệu chung về tính năng dự án}

Hiện nay, Amazon Web Services (AWS) là một trong những nền tảng điện toán đám mây phổ biến nhất phục vụ học tập, nghiên cứu và triển khai ứng dụng. Tuy nhiên, việc khởi tạo và quản trị tài nguyên thông qua giao diện web (AWS Management Console) đòi hỏi nhiều thao tác thủ công, điều hướng qua nhiều màn hình cấu hình phức tạp và rất dễ nhầm lẫn. Trong khi đó, việc sử dụng giao diện dòng lệnh (AWS CLI) hoặc viết tập lệnh tự động lại yêu cầu ghi nhớ chính xác cú pháp câu lệnh, cấu trúc tệp cấu hình JSON và các tham số kỹ thuật. Đối với sinh viên, lập trình viên mới tiếp cận điện toán đám mây hoặc người quản lý kỹ thuật, các rào cản này làm chậm quá trình vận hành, đồng thời tiềm ẩn nguy cơ cấu hình sai gây mất an toàn thông tin hoặc phát sinh chi phí ngoài ý muốn.

Nhằm giải quyết vấn đề trên, nhóm thực hiện đề tài xây dựng SatiBot -- ứng dụng ChatOps hỗ trợ quản lý, cấu hình và giám sát tài nguyên AWS thông qua câu lệnh bằng ngôn ngữ tự nhiên tiếng Việt. Bằng cách kết hợp mô hình ngôn ngữ lớn (LLM) hỗ trợ cơ chế \textit{Function Calling} với thư viện AWS SDK (Boto3), SatiBot hoạt động như một trợ lý hỗ trợ thao tác hạ tầng. Người dùng chỉ cần gửi yêu cầu bằng tiếng Việt qua giao diện chat; hệ thống sẽ phân tích ý định, kiểm tra tính hợp lệ và độ an toàn của câu lệnh, sau đó gọi các API tương ứng của AWS và trả về kết quả thực thi kèm theo thông tin chi tiết, cảnh báo rủi ro hoặc hướng dẫn xử lý tiếp theo.

\subsection{Giới thiệu người dùng của dự án}

Hệ thống SatiBot hướng tới 03 nhóm người dùng chính với nhu cầu sử dụng thực tế như sau:

\begin{itemize}
  \item \textbf{Lập trình viên và kỹ sư vận hành (Developers / DevOps Engineers):} Cần khởi tạo nhanh môi trường thử nghiệm, cấu hình máy chủ, cập nhật mã nguồn hàm serverless hoặc xem log hệ thống trong quá trình phát triển ứng dụng. SatiBot giúp nhóm này thực hiện các tác vụ lặp đi lặp lại nhanh hơn ngay trên khung chat mà không cần chuyển đổi qua lại giữa nhiều trang trên AWS Console.
  \item \textbf{Sinh viên ngành Công nghệ thông tin / Điện toán đám mây:} Cần môi trường thực hành trực quan cho các môn học về mạng và hạ tầng đám mây. Hệ thống cung cấp cơ chế kiểm tra an toàn (guardrails), giúp sinh viên tránh các lỗi phổ biến như mở cổng mạng công khai nguy hiểm, quên tắt máy ảo hoặc chọn nhầm các gói cấu hình vượt quá phạm vi gói miễn phí (AWS Free Tier).
  \item \textbf{Quản lý kỹ thuật và chủ dự án (Tech Leads / Project Managers):} Cần theo dõi nhanh tình trạng hoạt động của tài nguyên, kiểm tra phân quyền (IAM Audit) và kiểm soát chi phí sử dụng hàng tháng (FinOps) bằng ngôn ngữ tự nhiên mà không cần thành thạo cú pháp dòng lệnh AWS CLI.
\end{itemize}

\subsection{Cách thức và quy trình người dùng sử dụng sản phẩm để đạt mục đích}

\begingroup
\setlist{topsep=2pt,itemsep=1pt,parsep=0pt,leftmargin=*}
\setlength{\parskip}{2pt}

Để đáp ứng nhu cầu quản trị hạ tầng cơ bản, SatiBot hỗ trợ 04 dịch vụ chính của AWS gồm: Amazon EC2, Amazon S3, AWS IAM và AWS Lambda. Dưới đây là chi tiết về mục đích và quy trình tương tác thực tế của người dùng đối với từng nhóm tính năng của 04 dịch vụ AWS nòng cốt.

\subsubsection{Nhóm kịch bản và quy trình quản trị Amazon EC2}

Dịch vụ máy chủ ảo (Amazon Elastic Compute Cloud - EC2) là hạ tầng phục vụ việc triển khai và chạy ứng dụng. SatiBot hỗ trợ 05 quy trình quản lý trạng thái và cấu hình máy ảo như sau:

\begin{enumerate}
  \item \textbf{Khởi tạo máy chủ chuẩn hóa và xử lý trùng lặp tên:}
  \begin{itemize}
    \item \textit{Mục đích:} Tạo nhanh máy ảo để thử nghiệm hoặc triển khai website mà không phải cấu hình thủ công nhiều bước, đồng thời xử lý trường hợp tên máy bị trùng.
    \item \textit{Quy trình người dùng:} 
    \begin{itemize}
      \item Nhắn yêu cầu tạo máy kèm mục đích sử dụng (ví dụ: "Tạo cho tôi một máy chủ Ubuntu để chạy web demo, mở cổng SSH và HTTP").
      \item Nếu tên máy đã tồn tại, hệ thống đưa ra 3 lựa chọn: tiếp tục sử dụng máy cũ, đổi tên máy mới hoặc xóa máy cũ để tạo lại.
      \item Nhận kết quả thành công gồm: Địa chỉ IP máy chủ, link tải file khóa SSH (.pem) và câu lệnh kết nối mẫu.
    \end{itemize}
  \end{itemize}

  \item \textbf{Kiểm soát chi phí và xử lý vượt hạn mức tài nguyên (Cost Guardrails \& Quota):}
  \begin{itemize}
    \item \textit{Mục đích:} Tránh việc chọn nhầm máy cấu hình cao gây phát sinh chi phí lớn, đồng thời hướng dẫn cách xử lý khi tài khoản chạm ngưỡng hạn mức tạo máy ảo (vCPU quota).
    \item \textit{Quy trình người dùng:}
    \begin{itemize}
      \item Nhắn yêu cầu loại máy mong muốn (ví dụ: \textit{''Tạo máy ảo c5.4xlarge 16 vCPU''} hoặc \textit{''Tạo thêm 1 máy ảo Ubuntu''}).
      \item Xem phản hồi từ hệ thống:
      \begin{itemize}
        \item \textit{Nếu yêu cầu máy cấu hình cao, chi phí lớn:} Hệ thống cảnh báo chi phí ước tính hàng tháng, tạm dừng thao tác và yêu cầu người dùng xác nhận nếu vẫn muốn tiếp tục khởi tạo.
        \item \textit{Nếu tài khoản chạm giới hạn vCPU:} Hệ thống giải thích nguyên nhân và liệt kê danh sách các máy ảo đang chạy để người dùng cân nhắc tắt bớt máy không dùng.
      \end{itemize}
      \item Người dùng chọn phương án: chuyển về gói miễn phí (Free Tier), tắt bớt máy cũ hoặc gửi yêu cầu nâng hạn mức lên AWS.
    \end{itemize}
  \end{itemize}

  \item \textbf{Tra cứu và điều khiển vòng đời máy chủ (Bật / Tắt / Khởi động lại / Xóa):}
  \begin{itemize}
    \item \textit{Mục đích:} Theo dõi danh sách máy ảo và điều khiển trạng thái bật/tắt kịp thời nhằm tránh lãng phí tài nguyên và kiểm soát chi phí vận hành.
    \item \textit{Quy trình người dùng:}
    \begin{itemize}
      \item Nhắn câu lệnh tra cứu hoặc điều khiển (ví dụ: \textit{''Liệt kê các máy ảo đang chạy''}, \textit{''Tắt máy ảo web-demo''}, \textit{''Khởi động lại máy i-0abc123''}, \textit{''Xóa máy ảo staging''}).
      \item Xem bảng thông tin gồm ID, tên máy, trạng thái, IP công khai và thời gian chạy; nhận thông báo trạng thái cập nhật (kèm IP mới khi bật lại máy).
      \item Với thao tác xóa: Hệ thống cảnh báo mất dữ liệu và yêu cầu người dùng gõ xác nhận trước khi thực hiện xóa máy.
    \end{itemize}
  \end{itemize}

  \item \textbf{Giám sát hiệu năng và phát hiện máy chủ lãng phí (FinOps \& CloudWatch):}
  \begin{itemize}
    \item \textit{Mục đích:} Theo dõi mức độ sử dụng tài nguyên để tắt các máy nhàn rỗi nhằm tiết kiệm chi phí, hoặc nhận biết máy bị quá tải để kịp thời nâng cấp.
    \item \textit{Quy trình người dùng:}
    \begin{itemize}
      \item Nhắn câu hỏi kiểm tra (ví dụ: \textit{''Kiểm tra CPU máy i-0abc123 trong 24 giờ qua''}).
      \item Nhận số liệu đánh giá chi tiết:
      \begin{itemize}
        \item \textit{Cảnh báo lãng phí:} Nếu máy chạy không tải kéo dài, hệ thống đưa ra gợi ý tắt máy để tiết kiệm chi phí.
        \item \textit{Cảnh báo nghẽn tải:} Nếu mức sử dụng CPU liên tục ở mức cao, hệ thống cảnh báo quá tải và gợi ý nâng cấp loại máy.
      \end{itemize}
      \item Người dùng nhắn đồng ý nếu muốn hệ thống tự động xử lý theo khuyến nghị.
    \end{itemize}
  \end{itemize}

  \item \textbf{Thiết lập tường lửa an toàn và hỗ trợ sự cố kết nối (SecOps):}
  \begin{itemize}
    \item \textit{Mục đích:} Bảo vệ máy chủ khỏi nguy cơ bị quét dò mật khẩu từ Internet (brute-force attack) và hỗ trợ chẩn đoán nhanh nguyên nhân khi không kết nối SSH được vào máy ảo.
    \item \textit{Quy trình người dùng:}
    \begin{itemize}
      \item Người dùng gửi yêu cầu mở cổng hoặc hỏi nguyên nhân sự cố (ví dụ: \textit{''Mở port 22 cho mọi người truy cập''} hoặc \textit{''Tôi không SSH được vào máy ảo i-0abc123, sao vậy?''}).
      \item \textit{Kiểm tra an toàn tường lửa:} Hệ thống từ chối mở port 22 cho toàn bộ dải mạng công khai (0.0.0.0/0), giải thích nguy cơ bị tấn công dò quét và gợi ý chỉ mở riêng cho địa chỉ IP hiện tại của người dùng.
      \item \textit{Chẩn đoán kết nối:} Nếu phát hiện người dùng mất file khóa SSH (.pem), hệ thống giải thích nguyên nhân và hướng dẫn sử dụng các kênh truy cập thay thế (như AWS Systems Manager Session Manager hoặc EC2 Instance Connect).
      \item Người dùng phản hồi đồng ý để áp dụng cấu hình an toàn được đề xuất.
    \end{itemize}
  \end{itemize}
\end{enumerate}

\subsubsection{Nhóm kịch bản và quy trình quản trị Amazon S3}

Dịch vụ lưu trữ đối tượng (Amazon Simple Storage Service - S3) dùng để lưu trữ tập tin, bản sao lưu và dữ liệu tĩnh của hệ thống. SatiBot hỗ trợ 06 quy trình quản lý S3 như sau:

\begin{enumerate}
  \item \textbf{Khởi tạo thùng chứa (S3 Bucket) an toàn và chuẩn hóa:}
  \begin{itemize}
    \item \textit{Mục đích:} Tạo kho lưu trữ đối tượng được thiết lập ở chế độ riêng tư theo mặc định và tuân thủ quy tắc đặt tên của AWS.
    \item \textit{Quy trình người dùng:} Nhắn yêu cầu tạo kèm tên mong muốn (ví dụ: \textit{''Tạo bucket tailieu-backup''}). Nếu tên bị trùng hoặc sai quy chuẩn đặt tên, xem hướng dẫn từ hệ thống và nhắn lại tên mới. Nhận thông báo bucket đã được tạo thành công ở chế độ riêng tư.
  \end{itemize}

  \item \textbf{Quét và khóa rủi ro lộ dữ liệu ra Internet (Security Scan \& Remediation):}
  \begin{itemize}
    \item \textit{Mục đích:} Rà soát và ngăn chặn kịp thời tình trạng lộ dữ liệu do vô tình mở quyền truy cập công khai (Public Access) trên bucket.
    \item \textit{Quy trình người dùng:} Nhắn yêu cầu kiểm tra (ví dụ: \textit{''Kiểm tra xem có bucket nào đang bị lộ ra ngoài Internet không?''}). Hệ thống quét chính sách truy cập của toàn bộ bucket trong tài khoản và trả về danh sách các bucket đang có nguy cơ. Người dùng xác nhận (ví dụ: \textit{''Khóa lại giùm tôi''}) để hệ thống bật tính năng Block Public Access ngay lập tức.
  \end{itemize}

  \item \textbf{Xóa bucket S3 an toàn có kiểm soát (chống xóa nhầm dữ liệu):}
  \begin{itemize}
    \item \textit{Mục đích:} Dọn dẹp bucket cũ mà không lo ngại nguy cơ xóa nhầm dữ liệu quan trọng đang lưu bên trong.
    \item \textit{Quy trình người dùng:} Người dùng nhắn yêu cầu xóa (ví dụ: \textit{''Xóa bucket du-an-2025''}). Hệ thống kiểm tra và thông báo số lượng tệp tin còn tồn bên trong. Người dùng phải nhập xác nhận rõ ràng; khi đó hệ thống mới tiến hành dọn sạch các đối tượng và xóa bucket.
  \end{itemize}

  \item \textbf{Tối ưu chi phí lưu trữ lâu dài (FinOps Lifecycle Rules):}
  \begin{itemize}
    \item \textit{Mục đích:} Giảm thiểu chi phí lưu trữ bằng cách tự động chuyển các tệp cũ, ít sử dụng sang các tầng lưu trữ giá rẻ hơn (S3 Standard-IA, S3 Glacier).
    \item \textit{Quy trình người dùng:} Người dùng nhắn yêu cầu chuyển tầng lưu trữ kèm mốc thời gian (ví dụ: \textit{''Tự chuyển file trong bucket sang kho lưu trữ giá rẻ sau 30 ngày''}). Hệ thống thiết lập quy tắc vòng đời (Lifecycle Configuration Rule) trên S3 và gửi thông báo xác nhận chính sách đã được áp dụng.
  \end{itemize}

  \item \textbf{Chia sẻ tệp tin an toàn có thời hạn (Presigned URL):}
  \begin{itemize}
    \item \textit{Mục đích:} Cung cấp liên kết tải tệp tin an toàn cho người ngoài trong khoảng thời gian nhất định mà không cần mở public dữ liệu hay cấp quyền tài khoản AWS.
    \item \textit{Quy trình người dùng:} Người dùng nhắn tên tệp và thời hạn chia sẻ (ví dụ: \textit{''Tạo link tải file hop-dong-2026.pdf trong 15 phút''}). Hệ thống tạo đường dẫn có chữ ký xác thực (Presigned URL) với thời gian hết hạn (TTL) 15 phút và gửi liên kết trực tiếp trên khung chat.
  \end{itemize}

  \item \textbf{Tra cứu dung lượng và tình trạng kho lưu trữ (Storage Analytics):}
  \begin{itemize}
    \item \textit{Mục đích:} Nắm bắt nhanh số lượng tệp tin, dung lượng sử dụng thực tế và dự toán chi phí lưu trữ của từng bucket.
    \item \textit{Quy trình người dùng:} Người dùng đặt câu hỏi tra cứu (ví dụ: \textit{''Bucket tailieu-duan có bao nhiêu file và chiếm bao nhiêu GB?''}). Hệ thống truy vấn CloudWatch Metrics và S3 API để thống kê chi tiết: số lượng objects, dung lượng tổng (GB), tầng lưu trữ và chi phí ước tính, hiển thị dưới dạng bảng rõ ràng.
  \end{itemize}
\end{enumerate}

\subsubsection{Nhóm kịch bản và quy trình quản trị AWS IAM}

Dịch vụ quản lý định danh và quyền truy cập (AWS Identity and Access Management - IAM) là trung tâm kiểm soát an ninh của hạ tầng AWS. SatiBot cung cấp 05 quy trình quản trị IAM an toàn, tuân thủ nguyên tắc đặc quyền tối thiểu (Least Privilege):

\begin{enumerate}
  \item \textbf{Tra cứu vai trò (Role) và chính sách quyền hạn (Policy):}
  \begin{itemize}
    \item \textit{Mục đích:} Xem nhanh danh sách hoặc kiểm tra chi tiết các quyền hạn đang được gán cho nhân sự hoặc dịch vụ để phục vụ công tác kiểm toán an ninh.
    \item \textit{Quy trình người dùng:} Người dùng nhắn câu hỏi tra cứu (ví dụ: \textit{''Kiểm tra xem Role lambda-s3-role đang có những quyền gì?''} hoặc \textit{''Liệt kê các Role dành cho dịch vụ EC2''}). Hệ thống phân tích tài liệu chính sách JSON và hiển thị bảng tóm tắt quyền hạn (tên quyền, dịch vụ được phép truy cập, phạm vi đọc/ghi) ngay trên khung chat.
  \end{itemize}

  \item \textbf{Khởi tạo vai trò (Role) hoặc chính sách (Policy) mới theo nguyên tắc tối thiểu:}
  \begin{itemize}
    \item \textit{Mục đích:} Tạo mới đối tượng phân quyền cho máy chủ, ứng dụng hoặc nhân sự mà không phải tự tay viết các đoạn mã JSON chính sách phức tạp và dễ nhầm lẫn.
    \item \textit{Quy trình người dùng:} Người dùng nhắn yêu cầu tạo kèm mục đích (ví dụ: \textit{''Tạo một Role cho máy ảo EC2 chỉ được phép đọc dữ liệu từ S3''} hoặc \textit{''Tạo Policy cho phép đọc log CloudWatch''}). Hệ thống đề xuất cấu trúc quyền hạn phù hợp (chỉ cấp quyền đọc dữ liệu, không cấp thừa quyền ghi/xóa). Sau khi người dùng xác nhận, hệ thống gọi API tạo đối tượng và trả về ARN của Role/Policy.
  \end{itemize}

  \item \textbf{Cấp hoặc điều chỉnh quyền truy cập:}
  \begin{itemize}
    \item \textit{Mục đích:} Bổ sung hoặc nới rộng quyền hạn cho nhân sự/dịch vụ khi phát sinh công việc mới trong dự án.
    \item \textit{Quy trình người dùng:} Người dùng nhắn yêu cầu gán thêm quyền (ví dụ: \textit{''Cấp thêm quyền đọc DynamoDB cho Role app-backend''}). Hệ thống đối soát đối tượng nhận quyền và chính sách tương ứng, yêu cầu người dùng xác nhận và trả về thông báo quyền truy cập mới đã được áp dụng thành công.
  \end{itemize}

  \item \textbf{Thu hồi quyền truy cập (Least Privilege Enforcement):}
  \begin{itemize}
    \item \textit{Mục đích:} Gỡ bỏ các quyền hạn không còn sử dụng hoặc thu hồi các quyền nguy hiểm (như quyền AdministratorAccess cấp thừa cho tài khoản thực tập sinh) nhằm hạn chế rủi ro an toàn thông tin.
    \item \textit{Quy trình người dùng:} Người dùng nhắn yêu cầu gỡ bỏ quyền (ví dụ: \textit{''Gỡ quyền AdministratorAccess khỏi tài khoản dev-intern''}). Hệ thống hiển thị thông báo phạm vi quyền chuẩn bị gỡ bỏ; sau khi người dùng xác nhận, hệ thống gỡ bỏ chính sách và thông báo kết quả.
  \end{itemize}

  \item \textbf{Xóa đối tượng IAM an toàn:}
  \begin{itemize}
    \item \textit{Mục đích:} Dọn dẹp dứt điểm các Role, Policy hoặc User cũ không còn sử dụng để tránh dư thừa tài nguyên và loại bỏ lỗ hổng bảo mật tiềm tàng.
    \item \textit{Quy trình người dùng:} Người dùng gửi yêu cầu xóa (ví dụ: \textit{''Xóa Role test-migration cũ giùm tôi''}). Hệ thống kiểm tra các quan hệ phụ thuộc (nếu Role đang gắn với máy ảo EC2 đang chạy, hệ thống sẽ cảnh báo nguy cơ làm dừng ứng dụng). Người dùng xác nhận lần cuối để hệ thống tự động tháo gỡ các liên kết an toàn trước khi xóa bỏ đối tượng IAM.
  \end{itemize}
\end{enumerate}

\subsubsection{Nhóm kịch bản và quy trình quản trị AWS Lambda}

Dịch vụ điện toán không máy chủ (AWS Lambda) cho phép chạy mã nguồn theo sự kiện. SatiBot hỗ trợ 05 quy trình quản trị và kiểm thử hàm serverless trực tiếp:

\begin{enumerate}
  \item \textbf{Tra cứu danh mục và kiểm tra cấu hình hàm (Inventory \& Config Audit):}
  \begin{itemize}
    \item \textit{Mục đích:} Nắm bắt nhanh toàn bộ các hàm Lambda đang có trong tài khoản cùng các thông số vận hành cốt lõi (ngôn ngữ runtime, bộ nhớ RAM, thời gian timeout) phục vụ công tác kiểm toán tài nguyên.
    \item \textit{Quy trình người dùng:} Người dùng gửi câu hỏi (ví dụ: \textit{''Liệt kê các hàm Lambda trong tài khoản''} hoặc \textit{''Hàm image-processor đang chạy ngôn ngữ gì và cấp bao nhiêu RAM?''}). Hệ thống hiển thị bảng tổng hợp gồm Tên hàm, Môi trường thực thi (Runtime), Dung lượng bộ nhớ (Memory) và Thời gian chờ tối đa (Timeout).
  \end{itemize}

  \item \textbf{Chạy thử nghiệm hàm trực tiếp qua Chat (Direct Invocation):}
  \begin{itemize}
    \item \textit{Mục đích:} Kích hoạt và kiểm tra nhanh kết quả trả về của hàm với dữ liệu thử nghiệm ngay trên khung chat mà không cần đợi sự kiện bên ngoài hay cấu hình test event trên Console.
    \item \textit{Quy trình người dùng:} Người dùng nhắn lệnh chạy kèm dữ liệu đầu vào (ví dụ: \textit{''Chạy thử hàm calculate-price với dữ liệu giá 100 và thuế 10''}). Hệ thống đóng gói dữ liệu JSON, gọi API `invoke`, đo thời gian thực thi và trả về kết quả (Payload phản hồi, mã HTTP status, thời gian chạy tính bằng mili-giây) trực tiếp trên khung chat.
  \end{itemize}

  \item \textbf{Điều chỉnh cấu hình bộ nhớ và thời gian chờ của hàm Lambda:}
  \begin{itemize}
    \item \textit{Mục đích:} Nâng cấp cấu hình tức thời cho các hàm đang chạy chậm hoặc có nguy cơ sập do thiếu bộ nhớ hay hết thời gian chờ mà không phải can thiệp mã nguồn hay mở Console.
    \item \textit{Quy trình người dùng:} Người dùng ra lệnh điều chỉnh thông số (ví dụ: \textit{''Tăng thời gian timeout của hàm backup-data lên 30 giây''} hoặc \textit{''Nâng RAM của hàm image-processor lên 512 MB''}). Hệ thống hiển thị bảng đối soát thông số cũ và thông số mới. Sau khi người dùng xác nhận, hệ thống gọi API cập nhật cấu hình và thông báo thành công.
  \end{itemize}

  \item \textbf{Cấp phát nhanh liên kết HTTP công khai (Function URL Generation):}
  \begin{itemize}
    \item \textit{Mục đích:} Tạo nhanh một đường dẫn Endpoint HTTPS công khai (Function URL) để người dùng hoặc webhook bên ngoài có thể gọi trực tiếp hàm Lambda mà không mất công cấu hình API Gateway phức tạp.
    \item \textit{Quy trình người dùng:} Người dùng nhắn: \textit{''Tạo link Function URL công khai cho hàm image-api''}. Hệ thống tự động cấu hình Function URL, thiết lập quyền truy cập và gửi lại đường link HTTPS dạng \url{https://<url-id>.lambda-url.<region>.on.aws/} trực tiếp trên khung chat để người dùng kiểm chứng kết quả thực thi của hàm.
  \end{itemize}

  \item \textbf{Tạm dừng hoạt động hàm khẩn cấp (Emergency Kill-Switch):}
  \begin{itemize}
    \item \textit{Mục đích:} Vô hiệu hóa ngay lập tức một hàm đang bị lặp lỗi logic gây tiêu hao ngân sách hoặc cần ngắt kết nối gấp để bảo trì hệ thống mà không cần xóa bỏ mã nguồn của hàm.
    \item \textit{Quy trình người dùng:} Người dùng nhắn lệnh ngắt khẩn cấp (ví dụ: \textit{''Khóa tạm thời hàm order-service để bảo trì hệ thống''}). Hệ thống cảnh báo về việc hàm sẽ từ chối mọi yêu cầu mới (trả về lỗi HTTP 429 Too Many Requests) và đề nghị người dùng xác nhận. Sau khi nhận xác nhận, hệ thống đặt mức thực thi đồng thời về 0 (\texttt{Reserved Concurrency = 0}), đưa hàm vào trạng thái tạm ngừng hoạt động an toàn. Người dùng có thể yêu cầu mở lại mức concurrency bất cứ lúc nào.
  \end{itemize}
\end{enumerate}
\endgroup

\clearpage

\section{GIẢI PHÁP THỰC HIỆN CHÍNH THỨC}

\subsection{Kiến trúc tính năng của hệ thống}

\begin{figure}[H]
\centering
\IfFileExists{So_do_kien_truc_tinh_nang.jpg}{
  \includegraphics[width=0.8\textwidth]{So_do_kien_truc_tinh_nang.jpg}
}{
  \IfFileExists{Report progress 2/So_do_kien_truc_tinh_nang.jpg}{
    \includegraphics[width=0.8\textwidth]{Report progress 2/So_do_kien_truc_tinh_nang.jpg}
  }{
    \IfFileExists{Report progress 1/So_do_kien_truc_tinh_nang.jpg}{
      \includegraphics[width=0.8\textwidth]{Report progress 1/So_do_kien_truc_tinh_nang.jpg}
    }{
      \fbox{\parbox[c][6.5cm][c]{0.88\textwidth}{\centering Sơ đồ kiến trúc tính năng hệ thống SatiBot}}
    }
  }
}
\caption{Sơ đồ kiến trúc tính năng và luồng xử lý nghiệp vụ của SatiBot}
\label{fig:kientruc_tinhnang}
\end{figure}

Các tính năng hệ thống cung cấp bao gồm:
\begin{enumerate}[itemsep=3pt,topsep=2pt]
  \item \textbf{Đăng ký và xác thực người dùng:} Amazon Cognito xử lý đăng ký, đăng nhập và cấp JWT để xác thực các yêu cầu nghiệp vụ. Sau khi đăng ký, Lambda tạo hồ sơ người dùng trong bảng \texttt{Users} của DynamoDB.

  \item \textbf{Kết nối tài khoản AWS:} SatiBot tạo \texttt{ExternalId}; người dùng dùng mã này và hướng dẫn của hệ thống để tạo IAM Role trong tài khoản AWS của mình, sau đó gửi \texttt{RoleArn} về SatiBot. Connector Lambda gọi STS \texttt{AssumeRole} để kiểm tra khả năng ủy quyền. Nếu thành công, hệ thống lưu thông tin kết nối và trạng thái trong bảng \texttt{AwsConnections}.


  \begin{figure}[H]
  \centering
  \IfFileExists{sequence3.jpg}{
    \includegraphics[width=0.78\textwidth]{sequence3.jpg}
  }{
    \IfFileExists{Report progress 2/sequence3.jpg}{
      \includegraphics[width=0.78\textwidth]{Report progress 2/sequence3.jpg}
    }{
      \IfFileExists{Report progress 1/sequence3.jpg}{
        \includegraphics[width=0.78\textwidth]{Report progress 1/sequence3.jpg}
      }{
        \fbox{\parbox[c][5cm][c]{0.88\textwidth}{\centering Sơ đồ tuần tự quy trình kết nối tài khoản AWS}}
      }
    }
  }
  \caption{Sơ đồ tuần tự quy trình kết nối và xác thực tài khoản AWS mục tiêu}
  \label{fig:sequence_connect}
  \end{figure}

  \item \textbf{Quản trị tài nguyên qua Chatbot:} Chatbot Lambda nhận câu lệnh, lấy ngữ cảnh hội thoại từ \texttt{SessionState} và gửi yêu cầu cùng Function Calling Schema đến mô hình LLM. Kết quả được lưu thành hành động chờ xác nhận trong \texttt{PendingAction}. Khi người dùng xác nhận, Chatbot Lambda chuyển ActionID vào SQS cho Dispatcher Lambda; Dispatcher Lambda dùng STS và Boto3 để gọi API của EC2, S3, IAM hoặc Lambda trên AWS của người dùng. Kết quả được lưu vào \texttt{ActionHistory} và gửi về giao diện qua AppSync.

  \begin{figure}[H]
  \centering
  \IfFileExists{sequence1.jpg}{
    \includegraphics[width=0.8\textwidth]{sequence1.jpg}
  }{
    \IfFileExists{Report progress 2/sequence1.jpg}{
      \includegraphics[width=0.8\textwidth]{Report progress 2/sequence1.jpg}
    }{
      \IfFileExists{Report progress 1/sequence1.jpg}{
        \includegraphics[width=0.8\textwidth]{Report progress 1/sequence1.jpg}
      }{
      }
    }
  }
  \caption{Sơ đồ tuần tự quy trình điều phối và thực thi tác vụ ChatOps}
  \label{fig:sequence_chatops}
  \end{figure}

  \item \textbf{Tra cứu lịch sử thao tác:} History Lambda truy vấn bảng \texttt{ActionHistory} theo người dùng và các điều kiện tra cứu, sau đó trả về thời gian, hành động, tài nguyên đích và trạng thái thực thi để hiển thị trên giao diện.
\end{enumerate}

\subsection{Kiến trúc triển khai của hệ thống}

SatiBot sử dụng các dịch vụ AWS được quản lý để triển khai giao diện, xác thực, xử lý nghiệp vụ và lưu trữ dữ liệu.

\begin{table}[H]
\centering
\caption{Các thành phần và dịch vụ trong kiến trúc triển khai hệ thống SatiBot}
\label{tab:kientruc}
\bangfont
\begin{tabularx}{\linewidth}{@{}>{\raggedright\arraybackslash}p{4.5cm}Y@{}}
\toprule
\textbf{Thành phần / Dịch vụ} & \textbf{Vai trò trong kiến trúc}\\
\midrule
AWS WAF & Tường lửa ứng dụng web, gắn vào CloudFront để lọc request. \\
Amazon Route 53 & Quản lý tên miền (DNS), điều hướng người dùng đến CloudFront. \\
React, Amazon S3 và CloudFront & Cung cấp giao diện ChatOps; S3 lưu các tệp tĩnh và CloudFront phân phối giao diện qua HTTPS.\\
Amazon Cognito & Quản lý đăng ký, đăng nhập và cấp JWT để xác thực người dùng.\\
AWS STS & Cấp quyền truy cập tạm thời (Temporary Credentials). \\
Amazon API Gateway & Tiếp nhận API từ frontend, xác thực JWT và định tuyến yêu cầu chat, xác nhận thao tác hoặc tra cứu lịch sử.\\
AWS Step Functions & Điều phối luồng thực thi sau khi người dùng xác nhận, gồm ChatBot, Dispatcher và Formatter Lambda.\\
AWS Lambda & Thực hiện các nghiệp vụ: Chatbot, Dispatcher, Formatter, Connector và History; Connector đảm nhiệm tạo ExternalId và kiểm tra kết nối AWS.\\
Mô hình LLM & Phân tích câu lệnh tiếng Việt và trả về ý định, tham số theo Function Calling Schema; không trực tiếp gọi API AWS.\\
AWS STS và Boto3 & STS cấp thông tin xác thực tạm thời cho tài khoản đích; Boto3 dùng thông tin này để gọi API tài nguyên.\\
Amazon DynamoDB & Lưu hồ sơ người dùng, kết nối AWS, lịch sử hội thoại, hành động chờ xác nhận và lịch sử thao tác.\\
AWS AppSync & Phát kết quả qua GraphQL Subscriptions để giao diện cập nhật mà không cần polling liên tục.\\
Amazon CloudWatch & Lưu log thực thi và hỗ trợ theo dõi, chẩn đoán lỗi.\\
CI/CD Pipeline (GitHub, CodePipeline, CodeBuild) & Tự động hóa quy trình build, test và deploy. \\
\bottomrule
\end{tabularx}
\end{table}

\subsubsection{Sơ đồ kiến trúc triển khai hệ thống}

\begin{figure}[H]
\centering
\IfFileExists{so_do_kien_truc.jpg}{
  \includegraphics[width=0.82\textwidth]{so_do_kien_truc.jpg}
}{
  \IfFileExists{Report progress 2/so_do_kien_truc.jpg}{
    \includegraphics[width=0.82\textwidth]{Report progress 2/so_do_kien_truc.jpg}
  }{
    \IfFileExists{Report progress 1/so_do_kien_truc.jpg}{
      \includegraphics[width=0.82\textwidth]{Report progress 1/so_do_kien_truc.jpg}
    }{
      \fbox{\parbox[c][6.5cm][c]{0.88\textwidth}{\centering Sơ đồ kiến trúc triển khai tổng thể hệ thống SatiBot}}
    }
  }
}
\caption{Sơ đồ kiến trúc triển khai tổng thể hệ thống SatiBot}
\label{fig:so_do_kien_truc}
\end{figure}


Luồng hoạt động và giao tiếp giữa các thành phần gồm bốn bước:
\begin{enumerate}[itemsep=3pt,topsep=2pt]
  \item \textbf{Tiếp nhận yêu cầu:} Người dùng truy cập ứng dụng web thông qua giao thức HTTPS. Yêu cầu được gửi đến AWS Route53 và trả về IP của CloudFront. Request sẽ đi qua AWS WAF để chặn các request xấu (như chặn SQL Injection, Cross-Site Scripting (XSS), hoặc giới hạn tần suất request để chống DDoS). Sau đó, Amazon CloudFront là nơi phân phối nội dung tĩnh (HTML, CSS, JS) từ Amazon S3. Các yêu cầu API động được CloudFront chuyển tiếp đến Amazon API Gateway. API Gateway sử dụng Amazon Cognito để xác thực JWT Token của người dùng

  \item \textbf{Điều hướng yêu cầu:} Sau khi xác thực thành công, API Gateway kích hoạt dịch vụ tương ứng với yêu cầu của người dùng gồm Lambda for History (truy vấn lịch sử) và AWS core Step Function (Chatbot) .

  \item \textbf{Thực hiện yêu cầu:} Bên trong Step Functions Workflow, các hàm AWS Lambda được triệu gọi tuần tự:
\begin{itemize}
    \item Chatbot Lambda gọi Model LLM để phân tích ý định của người dùng và trích xuất Tool Call. Trạng thái chờ xác nhận được lưu vào bảng \texttt{PendingAction} trong Amazon DynamoDB.
    \item Sau khi người dùng xác nhận, Dispatcher Lambda gọi AWS STS để thực hiện \texttt{AssumeRole} lấy quyền truy cập tạm thời. Sau đó, nó sử dụng AWS SDK (\texttt{Boto3}) để thực thi các lệnh quản lý tài nguyên (như EC2, S3, IAM) trên tài khoản AWS cá nhân của người dùng.
    \item Formatter Lambda gọi lại Model LLM để chuyển đổi kết quả kỹ thuật thô thành ngôn ngữ tự nhiên, đồng thời ghi log hành động vào bảng \texttt{ActionHistory} trong DynamoDB và log hệ thống vào CloudWatch.
\end{itemize}

  \item \textbf{Ghi nhận và trả kết quả:} Formatter Lambda publish thông báo real-time qua AWS AppSync. AppSync duy trì kết nối WebSocket persistent và đẩy thông báo trực tiếp về cho Client/Web
\end{enumerate}

\subsection{Đặc tả chi tiết cần thiết trong giải pháp của dự án}
\subsubsection{Đặc tả cấu hình dịch vụ Bảo mật và Điều hướng}

\textbf{1. Amazon Route 53}
\begin{itemize}[leftmargin=*]
    \item \textbf{Tên miền \& Subdomain:} 
    \begin{itemize}[label=--]
        \item \code{satibot.id.vn} (Giao diện Web).
        \item \code{api.satibot.id.vn} (API Gateway Endpoint).
    \end{itemize}
    
    \item \textbf{Loại bản ghi:} Alias Record (trỏ tới CloudFront Distribution).
    
    \item \textbf{Chứng chỉ SSL/TLS:} AWS Certificate Manager (ACM).
    \begin{itemize}[label=--]
        \item Loại chứng chỉ: Wildcard (\code{*.satibot.id.vn}).
        \item Region cấp phát: \code{us-east-1}.
    \end{itemize}
     
    
    \item \textbf{Chính sách định tuyến:} Failover Routing.
    
    \item \textbf{Cấu hình bản ghi Primary:} Trỏ đến CloudFront Distribution (Region: \code{ap-southeast-1}).
    
    \item \textbf{Cấu hình bản ghi Secondary:} Trỏ đến CloudFront Distribution (Region: \code{us-east-1}; Origin: S3 Bucket chứa tệp \code{error.html}).
    
    \item \textbf{Cấu hình Health Check (áp dụng cho bản ghi Primary):}
    \begin{itemize}[label=--]
        \item Protocol: HTTPS.
        \item Path: \code{/health}.
        \item Request Interval: 10 giây.
        \item Failure Threshold: 3 lần.
    \end{itemize}
    
    \item \textbf{TTL:} 60 giây.
\end{itemize}

\vspace{0.5cm}

\textbf{2. AWS WAF (Web Application Firewall)}
\begin{itemize}[leftmargin=*]
    \item \textbf{Vị trí tích hợp:} Gắn vào CloudFront Distribution.
    
    \item \textbf{Cấu hình Web ACL:}
    \begin{itemize}[label=--]
        \item \textbf{Managed Rule Groups:}
        \begin{itemize}[label=]
            \item \code{AWSManagedRulesCommonRuleSet}: Action = Block.
            \item \code{AWSManagedRulesSQLiRuleSet}: Action = Block.
        \end{itemize}
        
        \item \textbf{Rate-based Rule:}
        \begin{itemize}[label=]
            \item Threshold: 100 requests / 5 phút / 1 IP.
            \item Action: Block.
        \end{itemize}
        
        \item \textbf{IP Sets:} Không cấu hình.
    \end{itemize}
\end{itemize}

\subsubsection{Giao diện và giao thức tương tác giữa các components}
\begin{table}[H]
\centering
\caption{Đặc tả giao diện và giao thức tương tác giữa các components}
\label{tab:inter_service_interfaces}
\begingroup
\renewcommand{\arraystretch}{1.12}
\fontsize{8.5}{11.2}\selectfont
\begin{tabularx}{\linewidth}{@{}>{\raggedright\arraybackslash}p{4.1cm}>{\raggedright\arraybackslash}p{2.35cm}>{\raggedright\arraybackslash}p{3.05cm}Y@{}}
\toprule
\textbf{Nguồn $\rightarrow$ Đích} & \textbf{Giao thức} & \textbf{Xác thực} & \textbf{Dữ liệu trao đổi}\\
\midrule
Web Client $\rightarrow$ CloudFront $\rightarrow$ S3 & HTTPS & TLS & Tệp giao diện tĩnh\\
Web Client $\rightarrow$ API Gateway & HTTPS REST & Cognito JWT & Câu lệnh, session ID, mã hành động\\
API Gateway $\rightarrow$ Chatbot Lambda & Lambda Proxy & IAM & Nội dung chat và thông tin phiên\\
API Gateway $\rightarrow$ Step Functions & StartExecution & IAM Role & User ID và mã hành động đã xác nhận\\
Chatbot Lambda $\leftrightarrow$ LLM & HTTPS API & Theo nhà cung cấp LLM & Prompt, schema công cụ, tham số và kết quả phân tích\\
Step Functions $\rightarrow$ Lambda Tasks & AWS SDK Integration & Step Functions Execution Role & JSON đầu vào/đầu ra từng bước\\
Lambda $\leftrightarrow$ DynamoDB & Boto3 SDK & Lambda Execution Role & GetItem, PutItem, Query, UpdateItem\\
Dispatcher Lambda $\rightarrow$ AWS STS & Boto3 / HTTPS & Lambda Execution Role & AssumeRole với RoleArn và ExternalId\\
Dispatcher $\rightarrow$ Tài khoản AWS đích & Boto3 / HTTPS & Credential tạm thời, SigV4 & Yêu cầu API và kết quả thao tác AWS\\
Formatter Lambda $\rightarrow$ AppSync & GraphQL Mutation / HTTPS & IAM SigV4 & Nội dung kết quả cần phát tới frontend\\
Web Client $\leftrightarrow$ AppSync & GraphQL Subscription / WSS & Cognito User Pool & Sự kiện cập nhật trên khung chat\\
\bottomrule
\end{tabularx}
\endgroup
\end{table}

\subsubsection{Cơ sở dữ liệu hệ thống}
\begin{figure}[H]
\centering
\IfFileExists{so_do_erd.jpg}{
  \includegraphics[width=0.9\textwidth]{so_do_erd.jpg}
}{
  \IfFileExists{Report progress 2/so_do_erd.jpg}{
    \includegraphics[width=0.9\textwidth]{Report progress 2/so_do_erd.jpg}
  }{
    \IfFileExists{Report progress 1/so_do_erd.jpg}{
      \includegraphics[width=0.9\textwidth]{Report progress 1/so_do_erd.jpg}
    }{
      \fbox{\parbox[c][6.5cm][c]{0.88\textwidth}{\centering Sơ đồ thực thể liên kết ERD}}
    }
  }
}
\caption{ERD cơ sở dữ liệu hệ thống}
\label{fig:so_do_erd}
\end{figure}



\subsubsection{CI/CD}

Quy trình CI/CD của SatiBot bắt đầu khi lập trình viên đẩy mã nguồn lên GitHub và tạo Pull Request. Mã nguồn được các thành viên trong nhóm xem xét, chạy kiểm thử đơn vị và chỉ được hợp nhất vào nhánh \texttt{main} khi đáp ứng yêu cầu. AWS CodeConnections theo dõi nhánh \texttt{main} và kích hoạt AWS CodePipeline khi xuất hiện commit mới. Trong giai đoạn tích hợp liên tục, AWS CodeBuild cài đặt Python cùng các thư viện cần thiết, chạy toàn bộ bộ kiểm thử bằng \texttt{pytest}, sau đó thực hiện \texttt{cdk synth} với hai biến cấu hình \texttt{PRIMARY\_REGION=ap-southeast-1} và \texttt{SECONDARY\_REGION=us-east-1}. Nếu quá trình build hoặc kiểm thử thất bại, pipeline dừng và phát thông báo lỗi; nếu thành công, quy trình chuyển sang bước phê duyệt thủ công. Khi yêu cầu triển khai được chấp thuận, AWS CDK và AWS CloudFormation tiến hành tạo mới hoặc cập nhật các tài nguyên AWS theo cấu hình đã khai báo.

Trong giai đoạn triển khai liên tục, hạ tầng được chia thành ba nhóm tài nguyên. Primary Stack tại \texttt{ap-southeast-1} chứa website chính và backend của SatiBot, bao gồm CloudFront, S3, API Gateway, WAF, Lambda, Cognito, IAM và DynamoDB. Các phiên bản Lambda mới được AWS CodeDeploy triển khai theo chiến lược canary; CloudWatch Logs và Alarm theo dõi trạng thái để chuyển toàn bộ lưu lượng sang phiên bản mới khi không có lỗi hoặc tự động rollback khi phát hiện sự cố. Error Page Stack tại \texttt{us-east-1} chỉ sử dụng S3 và CloudFront để cung cấp một trang thông báo lỗi độc lập, không cần đồng bộ dữ liệu ứng dụng với region chính. DNS/Failover Stack sử dụng Amazon Route~53 Health Check và chính sách Failover Routing để theo dõi hệ thống: khi region chính hoạt động bình thường, người dùng được điều hướng đến website tại \texttt{ap-southeast-1}; khi region này gặp sự cố, Route~53 chuyển kết quả phân giải DNS sang trang thông báo lỗi tại \texttt{us-east-1}.

\begin{figure}[H]
\centering
\IfFileExists{cicd.jpg}{
  \includegraphics[width=0.95\textwidth,height=0.78\textheight,keepaspectratio]{cicd.jpg}
}{
  \fbox{\parbox[c][7cm][c]{0.92\textwidth}{\centering Sơ đồ quy trình CI/CD của hệ thống SatiBot}}
}
\caption{Quy trình tích hợp, kiểm thử và triển khai liên tục của hệ thống SatiBot}
\label{fig:cicd_pipeline}
\end{figure}

\subsubsection{Phân quyền cho các hàm Lambda}

SatiBot phân quyền theo nhiệm vụ của từng Lambda. Mỗi hàm backend được gắn một IAM role để AWS xác định hàm được phép thực hiện thao tác nào và trên tài nguyên nào. Việc cấp quyền riêng giúp giới hạn tác động khi một hàm gặp lỗi hoặc xử lý sai yêu cầu.

Bảng~\ref{tab:lambda_permissions} trình bày các quyền theo thiết kế cho năm Lambda backend và Lambda ứng dụng trong tài khoản người dùng. Các quyền đọc, ghi được giới hạn vào đúng bảng dữ liệu; quyền đọc khóa Gemini API chỉ được cấp cho những hàm sử dụng mô hình.

\begingroup
\setlength{\tabcolsep}{5pt}
\renewcommand{\arraystretch}{1.17}
\fontsize{10.5}{14}\selectfont
\begin{longtable}{@{}>{\raggedright\arraybackslash}p{2.25cm}>{\raggedright\arraybackslash}p{3.35cm}>{\raggedright\arraybackslash}p{9.05cm}@{}}
\caption{Nhiệm vụ và quyền IAM của các Lambda}\label{tab:lambda_permissions}\\
\toprule
\textbf{Lambda} & \textbf{Nhiệm vụ} & \textbf{IAM Action cần cấp và nơi sử dụng}\\
\midrule
\endfirsthead
\toprule
\textbf{Lambda} & \textbf{Nhiệm vụ} & \textbf{IAM Action cần cấp và nơi sử dụng}\\
\midrule
\endhead
\endfoot
\bottomrule
\endlastfoot
Chatbot & Xử lý câu chat; lưu phiên và hành động chờ xác nhận. &
\path{dynamodb:GetItem}: đọc SessionState, PendingAction và AwsConnections.\par
\path{dynamodb:PutItem}, \path{dynamodb:UpdateItem}: ghi SessionState và PendingAction.\par
\path{secretsmanager:GetSecretValue}: đọc khóa Gemini API.\\
\midrule
Dispatcher & Thực thi hành động đã duyệt trên tài khoản AWS đích. &
\path{dynamodb:GetItem}: đọc PendingAction và AwsConnections.\par
\path{dynamodb:UpdateItem}: cập nhật PendingAction.\par
\path{sts:AssumeRole}: nhận quyền tạm thời từ role trong tài khoản đích.\\
\midrule
Formatter & Tạo lời đáp; lưu và gửi kết quả. &
\path{dynamodb:PutItem}, \path{dynamodb:UpdateItem}: ghi ActionHistory.\par
\path{appsync:GraphQL}: phát kết quả qua AppSync.\par
\path{secretsmanager:GetSecretValue}: đọc khóa Gemini API.\\
\midrule
Connector & Tạo ExternalId; lưu và kiểm tra kết nối AWS. &
\path{dynamodb:GetItem}, \path{dynamodb:PutItem}, \path{dynamodb:UpdateItem}: quản lý AwsConnections.\par
\path{sts:AssumeRole}: kiểm tra role người dùng cung cấp.\\
\midrule
History & Liệt kê và xem chi tiết lịch sử. &
\path{dynamodb:Query}: liệt kê bản ghi trong ActionHistory.\par
\path{dynamodb:GetItem}: đọc một bản ghi trong ActionHistory.\\
\midrule
Lambda của người dùng & Chạy mã ứng dụng trong tài khoản AWS đích. &
Quyền phụ thuộc mã ứng dụng; ví dụ \path{s3:GetObject} khi hàm cần đọc tệp S3. Quyền quản lý hàm, như \path{lambda:InvokeFunction} hoặc \path{lambda:UpdateFunctionConfiguration}, nằm trong role đích mà Dispatcher đảm nhận.\\
\end{longtable}
\endgroup

Để thao tác trên tài khoản AWS của người dùng, Connector và Dispatcher sử dụng role do chủ tài khoản cho phép SatiBot truy cập. Connector kiểm tra kết nối; Dispatcher nhận quyền tạm thời qua STS để thực hiện các thao tác được role đó cho phép. Lambda ứng dụng của người dùng có quyền chạy riêng, tùy theo mã ứng dụng. Các Lambda ghi log được cấp thêm quyền ghi CloudWatch Logs để theo dõi lỗi.

\subsubsection{Tích hợp LLM qua Gemini API}
\label{sec:llm_integration}

Hệ thống sử dụng mô hình Gemini của Google qua Gemini API. Tại Chatbot Lambda, mô hình chuyển câu chat tiếng Việt thành yêu cầu gọi hàm quản trị AWS; tại Formatter Lambda, mô hình diễn giải kết quả thực thi. Việc thao tác tài nguyên do Dispatcher thực hiện bằng Boto3 theo luồng ở Hình~\ref{fig:sequence_chatops}.


\begin{figure}[H]
\centering
\begin{tikzpicture}[x=1cm,y=1cm,
box/.style={draw,rounded corners,align=center,text width=3.6cm,minimum height=1cm,font=\small},
lab/.style={font=\footnotesize,align=center,fill=white,inner sep=2pt}]
\node[box] (chat) at (0,0) {Chatbot Lambda\\Chuẩn bị yêu cầu};
\node[box] (llm) at (8,0) {Gemini API\\Chọn hàm và tham số};
\node[box] (check) at (0,-2.5) {Chatbot Lambda\\Kiểm tra, yêu cầu xác nhận};
\node[box] (dispatch) at (8,-2.5) {Dispatcher Lambda\\Boto3 gọi AWS};
\draw[->,thick] (chat.north east) -- ++(0,.6) -| node[lab,pos=.25,above] {Câu chat, ngữ cảnh, danh sách hàm} (llm.north);
\draw[->,thick] (llm.west) -- node[lab,above] {Tên hàm và tham số} (chat.east);
\draw[->,thick] (chat) -- (check);
\draw[->,thick] (check) -- node[lab,above] {Hành động đủ điều kiện\\thực thi} (dispatch);
\end{tikzpicture}
\caption{Tích hợp Gemini API với các Lambda xử lý yêu cầu}
\label{fig:llm_integration}
\end{figure}

Backend gọi Gemini API bằng HTTP POST qua HTTPS, sử dụng endpoint Chat Completions do Google cung cấp. Nội dung request và response có định dạng JSON; khóa API được gửi trong header Authorization. Endpoint và mã model là cấu hình backend, model cần hỗ trợ Function Calling. Theo phương án triển khai, khóa API được lưu trong Secrets Manager và chỉ các Lambda cần gọi LLM được quyền đọc.

\begin{table}[H]
\centering
\caption{Dữ liệu trao đổi giữa Chatbot Lambda và Gemini API}
\label{tab:llm_exchange}
\fontsize{11}{14}\selectfont
\begin{tabularx}{\linewidth}{@{}>{\raggedright\arraybackslash}p{2.4cm}>{\raggedright\arraybackslash}p{3.2cm}Y@{}}
\toprule
\textbf{Chiều trao đổi} & \textbf{Trường dữ liệu} & \textbf{Nội dung}\\
\midrule
Gửi tới Gemini & \texttt{model} & Mã mô hình Gemini được cấu hình cho hệ thống.\\
 & \texttt{messages} & Chỉ dẫn hệ thống, câu chat hiện tại và ngữ cảnh hội thoại cần thiết.\\
 & \texttt{tools} & Danh sách hàm được phép gọi: tên, mô tả và JSON Schema quy định tham số.\\
 & \texttt{tool\_choice} & Giá trị \texttt{auto}: cho phép mô hình trả lời hoặc đề xuất gọi hàm theo yêu cầu.\\
\midrule
Nhận từ Gemini & \texttt{content} & Nội dung trả lời hoặc câu hỏi bổ sung khi chưa đủ thông tin.\\
 & \texttt{tool\_calls} & Danh sách lời gọi hàm; mỗi lời gọi có tên hàm và các tham số.\\
\bottomrule
\end{tabularx}
\end{table}

Danh sách hàm được khai báo trước cho bốn module EC2, S3, IAM và Lambda. Mỗi khai báo mô tả một thao tác và các tham số bắt buộc, kiểu dữ liệu, giá trị cho phép. Khi nhận \texttt{tool\_calls}, Chatbot kiểm tra tên hàm và tham số theo khai báo, rồi ánh xạ sang thao tác tương ứng của Dispatcher.

Formatter gửi kết quả AWS cho LLM để tạo lời đáp dễ đọc, không cần cung cấp danh sách hàm quản trị. Trạng thái thực thi, định danh tài nguyên và số liệu được lấy từ kết quả AWS; nội dung do LLM diễn giải không thay thế các dữ liệu này. Nếu bước diễn giải lỗi, Formatter trả kết quả theo mẫu có sẵn.

Mỗi lần gọi LLM có thời gian chờ và số lần thử lại giới hạn. Phản hồi sai cấu trúc không được chuyển sang thực thi. Ngữ cảnh hội thoại được giới hạn số lượt và dung lượng; credential AWS không được gửi tới mô hình. Lỗi kết nối hoặc lỗi tạo lời đáp không làm chạy lại thao tác thay đổi tài nguyên đã thực hiện.


\subsubsection{Endpoint tiếp nhận hội thoại và trích xuất ý định (POST /chat)}

\textbf{1. Thông tin chung:}
\begin{itemize}
    \item \textbf{Đường dẫn (Path):} \texttt{/chat} (hoặc \texttt{/api/chat} trên giao diện Web Client).
    \item \textbf{Phương thức HTTP:} \texttt{POST}.
    \item \textbf{Cơ chế xác thực:} Bearer JWT Token thông qua HTTP Header \texttt{Authorization: Bearer <Token>}. Amazon API Gateway sử dụng Cognito User Pool Authorizer để giải mã và kiểm tra tính hợp lệ của token trước khi kích hoạt hàm xử lý phía sau.
    \item \textbf{Thành phần tiếp nhận:} Amazon API Gateway $\rightarrow$ Chatbot Lambda (tích hợp mô hình ngôn ngữ lớn LLM với cơ chế Function Calling).
\end{itemize}

\textbf{2. Dữ liệu trao đổi thực tế (Kịch bản khởi tạo máy chủ ảo EC2):}

Mẫu JSON Request gửi từ Web Client khi người dùng gửi câu lệnh khởi tạo máy chủ EC2 bằng ngôn ngữ tự nhiên tiếng Việt:

\begin{verbatim}
{
  "message": "Tạo cho tôi một máy chủ EC2 Ubuntu có tên 'web-demo'",
  "history": [
    {
      "role": "user",
      "content": "Xin chào SatiBot"
    },
    {
      "role": "assistant",
      "content": "Xin chào! Tôi là SatiBot - Trợ lý ChatOps quản trị AWS."
    }
  ],
  "region": "ap-southeast-2"
}
\end{verbatim}

Mẫu JSON Response trả về từ hệ thống. Do thao tác \texttt{create\_instance} là hành động làm thay đổi tài nguyên hạ tầng và phát sinh chi phí, hệ thống trả về mã trạng thái \texttt{CONFIRMATION\_REQUIRED} kèm mã giao dịch \texttt{action\_id} để yêu cầu xác nhận:

\begin{verbatim}
{
  "status": "CONFIRMATION_REQUIRED",
  "type": "action_confirmation",
  "action_id": "act_8f3a9e1b2c4d",
  "service": "ec2",
  "tool_name": "create_instance",
  "parameters": {
    "instance_name": "web-demo",
    "os_type": "ubuntu",
    "instance_type": "t3.micro"
  },
  "message": "Thao tác 'create_instance' làm thay đổi cấu hình hoặc kích hoạt"
             " tài nguyên. Bạn có chắc chắn muốn thực hiện không?"
}
\end{verbatim}

\textbf{3. Bảng giải thích chi tiết các trường dữ liệu:}

\begin{table}[H]
\centering
\caption{Mô tả các trường trong Request Body của Endpoint \texttt{POST /chat}}
\label{tab:chat_request}
\begingroup
\renewcommand{\arraystretch}{1.12}
\fontsize{10}{13}\selectfont
\begin{tabularx}{\linewidth}{@{}>{\raggedright\arraybackslash}p{2.5cm}>{\raggedright\arraybackslash}p{2.3cm}>{\centering\arraybackslash}p{1.8cm}Y@{}}
\toprule
\textbf{Tên trường} & \textbf{Kiểu dữ liệu} & \textbf{Bắt buộc} & \textbf{Mô tả ý nghĩa và nguồn gốc dữ liệu} \\
\midrule
\texttt{message} & String & Có & Nội dung câu lệnh bằng ngôn ngữ tự nhiên tiếng Việt của người dùng gửi từ giao diện. \\
\texttt{history} & Array[Object] & Không & Lịch sử hội thoại nhiều lượt gồm các cặp \texttt{role} và \texttt{content}, đồng bộ từ bảng \texttt{SessionState} trong DynamoDB. \\
\texttt{region} & String & Không & Khu vực AWS mục tiêu thực hiện tác vụ (mặc định: \texttt{ap-southeast-2}). \\
\texttt{credentials} & Object & Không & Thông tin xác thực AWS cục bộ. \\
\bottomrule
\end{tabularx}
\endgroup
\end{table}

\begin{table}[H]
\centering
\caption{Mô tả các trường trong Response Body của Endpoint \texttt{POST /chat}}
\label{tab:chat_response}
\begingroup
\renewcommand{\arraystretch}{1.12}
\fontsize{10}{13}\selectfont
\begin{tabularx}{\linewidth}{@{}>{\raggedright\arraybackslash}p{2.8cm}>{\raggedright\arraybackslash}p{2.2cm}>{\centering\arraybackslash}p{1.8cm}Y@{}}
\toprule
\textbf{Tên trường} & \textbf{Kiểu dữ liệu} & \textbf{Bắt buộc} & \textbf{Mô tả ý nghĩa và nguồn gốc dữ liệu} \\
\midrule
\texttt{status} & String & Có & Trạng thái phản hồi của hệ thống (\texttt{CONFIRMATION\_REQUIRED} khi cần duyệt, hoặc \texttt{SUCCESS} khi chỉ truy vấn dữ liệu). \\
\texttt{type} & String & Có & Loại phản hồi định hướng hiển thị phía Client (\texttt{action\_confirmation} hoặc \texttt{chat\_reply}). \\
\texttt{action\_id} & String & Có & Mã định danh duy nhất cho hành động chờ duyệt, được lưu làm Partition Key trong bảng \texttt{PendingAction} với thời gian TTL 15 phút. \\
\texttt{service} & String & Có & Dịch vụ AWS mục tiêu (ví dụ: \texttt{ec2}, \texttt{s3}, \texttt{lambda}, \texttt{iam}) do LLM Function Calling trích xuất. \\
\texttt{tool\_name} & String & Có & Tên hành động kỹ thuật được trích xuất (ví dụ: \texttt{create\_instance}, \texttt{manage\_power}). \\
\texttt{parameters} & Object & Có & Cấu trúc tham số kỹ thuật chi tiết trích xuất từ câu lệnh tự nhiên của người dùng. \\
\texttt{message} & String & Có & Lời nhắn tiếng Việt thân thiện gửi người dùng để thông báo hoặc yêu cầu xác nhận an toàn. \\
\bottomrule
\end{tabularx}
\endgroup
\end{table}

\subsubsection{Endpoint xác nhận và thực thi tài nguyên (POST /confirm)}

\textbf{1. Thông tin chung:}
\begin{itemize}
    \item \textbf{Đường dẫn (Path):} \texttt{/confirm} (hoặc gửi tới \texttt{/chat} với tham số \texttt{confirmed: true}).
    \item \textbf{Phương thức HTTP:} \texttt{POST}.
    \item \textbf{Cơ chế xác thực:} Bearer JWT Token của Amazon Cognito trong Header \texttt{Authorization}.
    \item \textbf{Thành phần tiếp nhận:} Amazon API Gateway $\rightarrow$ Chatbot Lambda / Step Functions $\rightarrow$ Dispatcher Lambda $\rightarrow$ EC2 Service (\texttt{boto3.client('ec2')}) $\rightarrow$ Amazon DynamoDB (\texttt{ActionHistory}).
\end{itemize}

\textbf{2. Dữ liệu trao đổi thực tế (Kịch bản người dùng chấp thuận tạo máy ảo):}

Mẫu JSON Request gửi từ Web Client khi người dùng nhấn nút ``Xác nhận thực thi'' trên giao diện:

\begin{verbatim}
{
  "confirmed": true,
  "action_id": "act_8f3a9e1b2c4d",
  "service": "ec2",
  "tool_name": "create_instance",
  "parameters": {
    "instance_name": "web-demo",
    "os_type": "ubuntu",
    "instance_type": "t3.micro"
  },
  "region": "ap-southeast-2"
}
\end{verbatim}

Mẫu JSON Response trả về ngay lập tức sau khi Dispatcher Lambda và Module EC2 Service thực thi gọi Boto3 SDK thành công trên AWS Sydney:

\begin{verbatim}
{
  "status": "SUCCESS",
  "type": "action_executed",
  "service": "ec2",
  "tool_name": "create_instance",
  "message": "Đã gửi yêu cầu khởi tạo máy chủ 'web-demo' (i-015aad95042137744)"
             " thành công với trạng thái 'pending'.",
  "data": {
    "instance_id": "i-015aad95042137744",
    "instance_name": "web-demo",
    "state": "pending",
    "instance_type": "t3.micro",
    "image_id": "ami-0af9110f0ec0820e2"
  }
}
\end{verbatim}

\textbf{3. Bảng giải thích chi tiết các trường dữ liệu:}

\begin{table}[H]
\centering
\caption{Mô tả các trường trong Request Body của Endpoint \texttt{POST /confirm}}
\label{tab:confirm_request}
\begingroup
\renewcommand{\arraystretch}{1.12}
\fontsize{10}{13}\selectfont
\begin{tabularx}{\linewidth}{@{}>{\raggedright\arraybackslash}p{2.5cm}>{\raggedright\arraybackslash}p{2.3cm}>{\centering\arraybackslash}p{1.8cm}Y@{}}
\toprule
\textbf{Tên trường} & \textbf{Kiểu dữ liệu} & \textbf{Bắt buộc} & \textbf{Mô tả ý nghĩa và nguồn gốc dữ liệu} \\
\midrule
\texttt{confirmed} & Boolean & Có & Cờ xác nhận mang giá trị \texttt{true}, khẳng định người dùng đã trực tiếp phê duyệt hành động. \\
\texttt{action\_id} & String & Có & Mã hành động tương ứng đã tạo ở bước \texttt{/chat}, dùng để đối soát trạng thái trong bảng \texttt{PendingAction}. \\
\texttt{service} & String & Có & Dịch vụ AWS đích cần điều phối thực thi (\texttt{ec2}, \texttt{s3}, \texttt{lambda}, \texttt{iam}). \\
\texttt{tool\_name} & String & Có & Tên hành động kỹ thuật cần thực thi (\texttt{create\_instance}, v.v.). \\
\texttt{parameters} & Object & Có & Tập tham số kỹ thuật đã qua rà soát của người quản trị. \\
\texttt{region} & String & Không & Khu vực thực thi lệnh trên đám mây AWS (mặc định: \texttt{ap-southeast-2}). \\
\bottomrule
\end{tabularx}
\endgroup
\end{table}

\begin{table}[H]
\centering
\caption{Mô tả các trường trong Response Body của Endpoint \texttt{POST /confirm}}
\label{tab:confirm_response}
\begingroup
\renewcommand{\arraystretch}{1.12}
\fontsize{10}{13}\selectfont
\begin{tabularx}{\linewidth}{@{}>{\raggedright\arraybackslash}p{3.2cm}>{\raggedright\arraybackslash}p{2.0cm}>{\centering\arraybackslash}p{1.8cm}Y@{}}
\toprule
\textbf{Tên trường} & \textbf{Kiểu dữ liệu} & \textbf{Bắt buộc} & \textbf{Mô tả ý nghĩa và nguồn gốc dữ liệu} \\
\midrule
\texttt{status} & String & Có & Kết quả xử lý hạ tầng tổng thể (\texttt{SUCCESS} hoặc \texttt{ERROR}). \\
\texttt{type} & String & Có & Phân loại giao diện (\texttt{action\_executed}) để Web Client cập nhật trạng thái thẻ hiển thị sang hoàn tất. \\
\texttt{service} & String & Có & Tên dịch vụ AWS đã thực thi hoàn tất (\texttt{ec2}). \\
\texttt{tool\_name} & String & Có & Tên công cụ kỹ thuật đã chạy (\texttt{create\_instance}). \\
\texttt{message} & String & Có & Thông báo tổng kết kỹ thuật bằng tiếng Việt kèm mã ID và trạng thái máy ảo. \\
\texttt{data} & Object & Có & Đối tượng chứa dữ liệu chi tiết trả về từ Boto3 SDK gồm: \texttt{instance\_id}, \texttt{instance\_name}, \texttt{state} (trạng thái \texttt{pending} non-blocking), \texttt{instance\_type} và mã \texttt{image\_id} chính thức của hệ điều hành. \\
\bottomrule
\end{tabularx}
\endgroup
\end{table}


\section{KẾ HOẠCH TRIỂN KHAI DỰ ÁN}

\subsection{Kế hoạch triển khai theo từng giai đoạn (05 tuần)}

Dự án được lên kế hoạch thực hiện trong 5 tuần theo mô hình Agile/Scrum, chia thành 5 Sprint tương ứng:

\begin{itemize}
  \item \textbf{Tuần 1 (Sprint 1) -- Chuẩn bị môi trường, chốt thiết kế và kiểm thử PoC:} Thống nhất quy trình GitHub, phân công và tiêu chí nghiệm thu; chuẩn bị tài khoản AWS và khung hạ tầng CDK. Kiểm thử độc lập Boto3 cho S3, EC2, Lambda, IAM; chốt giao diện, hợp đồng API và thử nghiệm Gemini Function Calling.
  \item \textbf{Tuần 2 (Sprint 2) -- Dựng hạ tầng nền tảng và tích hợp Backend:} Triển khai React trên S3/CloudFront, Cognito và API Gateway; dựng khung năm Lambda Chatbot, Dispatcher, Formatter, Connector, History. Tạo 05 bảng DynamoDB, phân quyền riêng cho từng Lambda, lưu khóa Gemini trong Secrets Manager; triển khai kết nối AWS qua ExternalId/STS và pipeline build, test, synth hạ tầng.
  \item \textbf{Tuần 3 (Sprint 3) -- Hoàn thiện ChatOps, Gemini và cập nhật thời gian thực:} Hoàn thiện bốn module quản trị và Tool Schema; tích hợp Gemini tại Chatbot và Formatter. Xây dựng luồng xác nhận qua PendingAction và Step Functions, Dispatcher thực thi bằng STS/Boto3, Formatter lưu ActionHistory và phát kết quả qua AppSync; bổ sung tra cứu lịch sử trên giao diện.
  \item \textbf{Tuần 4 (Sprint 4) -- Bổ sung bảo vệ, dự phòng và kiểm thử toàn diện:} Cấu hình WAF, Route~53/ACM và chuyển hướng sang trang lỗi tại \texttt{us-east-1}. Hoàn thiện CDK/CloudFormation, phê duyệt triển khai, canary và rollback bằng CodeDeploy/CloudWatch. Kiểm thử nghiệp vụ, phân quyền, xác nhận, lỗi Gemini, kết nối thời gian thực và tối ưu thời gian phản hồi.
  \item \textbf{Tuần 5 (Sprint 5) -- Nghiệm thu, bàn giao và tổng kết:} Nghiệm thu toàn bộ kịch bản của bốn dịch vụ, luồng ChatOps, pipeline triển khai và failover; hoàn thiện hướng dẫn sử dụng, vận hành, báo cáo và slide. Rà soát, dọn tài nguyên thử nghiệm cùng các cấu hình phát sinh chi phí sau khi hoàn tất demo.
\end{itemize}

\clearpage
\subsection{Kế hoạch hành động chi tiết theo từng Sprint}

\begingroup
\setlength{\tabcolsep}{4pt}
\renewcommand{\arraystretch}{1.13}
\fontsize{11}{14.5}\selectfont
\begin{longtable}{@{}>{\centering\arraybackslash}p{0.75cm}>{\raggedright\arraybackslash}p{4.6cm}>{\raggedright\arraybackslash}p{3.2cm}>{\raggedright\arraybackslash}p{\dimexpr\linewidth-8.55cm-6\tabcolsep\relax}@{}}
\caption{Kế hoạch hành động và phân công công việc chi tiết theo từng Sprint}\label{tab:action_plan}\\
\toprule
\textbf{STT} & \textbf{Hạng mục công việc} & \textbf{Người phụ trách} & \textbf{Kết quả bàn giao}\\
\midrule
\endfirsthead
\toprule
\textbf{STT} & \textbf{Hạng mục công việc} & \textbf{Người phụ trách} & \textbf{Kết quả bàn giao}\\
\midrule
\endhead
\midrule
\endfoot
\bottomrule
\endlastfoot
\multicolumn{4}{@{}p{\linewidth}@{}}{\textbf{Giai đoạn 1: Tuần 1 -- Thiết lập quy trình, hạ tầng ban đầu và PoC Boto3 độc lập}}\\*
1 & Viết và kiểm thử độc lập module Boto3 cho dịch vụ S3 & \mbox{Phạm Đức Hùng} & Module \texttt{s3\_service} (tạo bucket, presign, block public).\\
2 & Thống nhất kịch bản và kiểm thử module Boto3 cho dịch vụ EC2 & \mbox{Nguyễn Đức Huy} & Module \texttt{ec2\_service} (khởi tạo, bật/tắt, describe).\\
3 & Thống nhất kịch bản và kiểm thử module Boto3 cho Lambda & \mbox{Trịnh Thiên Quang} & Module \texttt{lambda\_service} (list, invoke, hot config).\\
4 & Thống nhất quy tắc làm việc, tạo tài khoản AWS và module IAM\par Chuẩn bị khung hạ tầng AWS CDK cho mô hình mở rộng. & \mbox{Cao Anh Minh} & Quy chuẩn mã nguồn, AWS Account chung, module \texttt{iam\_service}.\par Khung Primary Stack, Error Page Stack và DNS/Failover Stack; cấu hình region chính và region trang lỗi.\\
\midrule
\multicolumn{4}{@{}p{\linewidth}@{}}{\textbf{Giai đoạn 2: Tuần 2 -- Dựng hạ tầng nền tảng Serverless, Giao diện và Tích hợp Backend}}\\*
5 & Phát triển khung giao diện React cơ bản trên S3/CloudFront\par Cấu hình AWS WAF cho CloudFront; dựng trang thông báo lỗi trên S3/CloudFront ở region phụ. & \mbox{Phạm Đức Hùng} & Giao diện Web Chat responsive, CDN CloudFront với HTTPS.\par Web ACL với managed rules và rate limit theo đặc tả; trang error.html tại us-east-1; script build frontend dùng trong CI/CD.\\
6 & Cấu hình luồng đăng nhập và quản lý danh tính qua Cognito\par Cấu hình tên miền web/API qua Route 53 và chứng chỉ ACM; chuẩn bị health check. & \mbox{Nguyễn Đức Huy} & Luồng Login/Register qua Cognito, JWT Authorizer trên Gateway.\par Bản ghi DNS, chứng chỉ cho tên miền web/API và cấu hình kiểm tra /health phục vụ failover.\\
7 & Dựng Lambda Backend trung tâm và tích hợp các module Boto3\par Mở rộng Backend theo các vai trò Chatbot, Dispatcher, Formatter, Connector, History; thiết lập CI. & \mbox{Trịnh Thiên Quang} & Lambda Backend tiếp nhận request và tích hợp các module Boto3.\par Khung năm Lambda, API kết nối AWS/tra cứu lịch sử; CodeConnections, CodePipeline, CodeBuild chạy pytest, build frontend và cdk synth.\\
8 & Cấu hình IAM Role cho Backend và thiết kế 05 bảng DynamoDB\par Cấp quyền riêng cho từng Lambda, Step Functions và pipeline; quản lý khóa Gemini bằng Secrets Manager. & \mbox{Cao Anh Minh} & IAM Execution Role, 05 bảng DynamoDB cơ sở.\par Policy theo chức năng, quyền đọc/ghi đúng bảng; chỉ Chatbot/Formatter đọc khóa Gemini; mẫu role đích có ExternalId và cấu hình CDK nền tảng.\\
\midrule
\multicolumn{4}{@{}p{\linewidth}@{}}{\textbf{Giai đoạn 3: Tuần 3 -- Tích hợp AI Chatbot, Step Functions và AppSync}}\\*
9 & Định nghĩa Function Calling và Prompt cho dịch vụ S3\par Ghép thẻ xác nhận, kết quả thực thi và lịch sử trên giao diện với Backend mở rộng. & \mbox{Phạm Đức Hùng} & Bộ Tool Schema cho S3, hỗ trợ kênh WebSocket AppSync.\par Giao diện nhận sự kiện AppSync theo người dùng/phiên; hiển thị kết quả và lịch sử từ API Backend.\\
10 & Định nghĩa Function Calling cho EC2 và ghép luồng LLM\par Hoàn thiện DNS failover giữa website chính và trang thông báo lỗi. & \mbox{Nguyễn Đức Huy} & Bộ Tool Schema cho EC2, hoàn thiện ghép luồng LLM.\par Health check và bản ghi Primary/Secondary; kết quả thử chuyển sang trang lỗi và phục hồi website chính.\\
11 & Tích hợp API LLM, xử lý phản hồi Function Calling tại Backend\par Tích hợp Gemini vào Chatbot/Formatter; hoàn thiện các Lambda nghiệp vụ của Backend mở rộng. & \mbox{Trịnh Thiên Quang} & Xử lý phản hồi Function Calling và tích hợp LLM vào Backend.\par Chatbot lưu phiên/hành động chờ; Dispatcher thực thi Boto3; Formatter lưu lịch sử, phát AppSync và có lời đáp dự phòng; Connector kiểm tra STS, History truy vấn lịch sử.\\
12 & Hoàn thiện Prompt IAM, xây dựng Step Functions State Machine\par Ghép State Machine với Backend mở rộng và luồng xác nhận hành động. & \mbox{Cao Anh Minh} & State Machine điều phối luồng thực thi, xác thực STS.\par Workflow gọi các Lambda đúng thứ tự; đối soát hành động, chủ sở hữu và hạn xác nhận; phối hợp Backend chống thực thi lặp.\\
\midrule
\multicolumn{4}{@{}p{\linewidth}@{}}{\textbf{Giai đoạn 4: Tuần 4 -- Kiểm thử toàn diện, Phân quyền và Tối ưu hóa}}\\*
13 & Kiểm thử hệ thống, tối ưu hóa module S3\par Kiểm thử WAF, trang lỗi và cập nhật giao diện qua AppSync. & \mbox{Phạm Đức Hùng} & Báo cáo kiểm thử S3, tối ưu giao diện truyền tải file.\par Kết quả kiểm thử request hợp lệ/bị chặn; hiển thị trang lỗi, kết nối lại AppSync và đối soát kết quả trên giao diện.\\
14 & Kiểm thử hệ thống, tối ưu module EC2, dọn dẹp tài nguyên\par Kiểm thử DNS/TLS, health check và failover. & \mbox{Nguyễn Đức Huy} & Báo cáo kiểm thử EC2, kịch bản FinOps cảnh báo máy ảo.\par Bằng chứng chuyển sang trang lỗi khi hệ thống chính lỗi và phục hồi khi hoạt động lại; kiểm tra tên miền/chứng chỉ web/API.\\
15 & Kiểm thử phân quyền, tối ưu hệ thống Backend\par Hoàn thiện và kiểm thử CI/CD, triển khai canary/rollback; kiểm thử các Lambda mở rộng. & \mbox{Trịnh Thiên Quang} & Báo cáo kiểm thử phân quyền, tối ưu Cold-start của Lambda.\par Pipeline phê duyệt thủ công, triển khai CDK/CloudFormation và CodeDeploy; kiểm thử lỗi Gemini, Connector, History và xác nhận lặp; bằng chứng canary/rollback.\\
16 & Kiểm thử trường hợp thiếu quyền (AccessDenied), CloudWatch Log\par Cấu hình CloudWatch Alarm cho CodeDeploy; rà soát quyền secret, workflow và pipeline. & \mbox{Cao Anh Minh} & Xử lý lỗi phân quyền, tập trung hóa log trên CloudWatch.\par Alarm phục vụ rollback; báo cáo AccessDenied, quyền giữa các Lambda và truy cập chéo người dùng; log truy vết cho Backend mở rộng.\\
\midrule
\multicolumn{4}{@{}p{\linewidth}@{}}{\textbf{Giai đoạn 5: Tuần 5 -- Nghiệm thu, Tổng kết dự án và Đóng gói báo cáo}}\\*
17 & Hoàn thiện báo cáo phần S3\par Bổ sung nội dung WAF, giao diện mở rộng và trang thông báo lỗi vào báo cáo/demo. & \mbox{Phạm Đức Hùng} & Nội dung báo cáo S3.\par Tài liệu cấu hình WAF, S3/CloudFront dự phòng; hình ảnh và kịch bản demo giao diện mới.\\
18 & Tổng hợp kết quả nghiệm thu EC2\par Tổng hợp nghiệm thu Route 53/ACM và failover; ghi nhận dọn tài nguyên thử nghiệm. & \mbox{Nguyễn Đức Huy} & Nội dung báo cáo EC2.\par Kết quả DNS/TLS, health check, chuyển hướng/phục hồi; danh sách tài nguyên EC2/DNS thử nghiệm đã xử lý.\\
19 & Nghiệm thu luồng Lambda, hoàn thiện tài liệu hướng dẫn\par Bổ sung tài liệu Backend mở rộng, Gemini và CI/CD. & \mbox{Trịnh Thiên Quang} & Nội dung báo cáo Lambda, tài liệu vận hành.\par Hướng dẫn năm Lambda, cấu hình Gemini, build/deploy, canary/rollback và xử lý lỗi; kết quả nghiệm thu pipeline.\\
20 & Tổng kết dự án, hoàn thiện báo cáo tổng thể\par Tổng hợp phần kiến trúc mở rộng và rà soát chi phí sau demo. & \mbox{Cao Anh Minh} & Bản báo cáo hoàn chỉnh.\par Nội dung tổng hợp ba stack, phân quyền, bảo mật và dự phòng; checklist dọn tài nguyên thử nghiệm và đối soát chi phí của cả nhóm.\\
\end{longtable}
\endgroup

\end{document}
